\chapter{tmux}\label{ch:tmux}

tmux runs many terminals inside one window and keeps them alive when the window closes. Config:
\file{tmux/.config/tmux/tmux.conf}; there are no plugins to install.

\section{Concepts}

\begin{keys}{2.6cm}
\head{Term}{Meaning}
Session & A named workspace, for example \texttt{main}, or one per project or per robot. Sessions keep running in the background.\\
Window & A tab inside a session, shown in the status bar as \texttt{1 molly}, \texttt{2 humble}.\\
Pane & A split inside a window.\\
Prefix & \key{C-a}: press it, release, then press the next key.\\
\end{keys}

Start or re-attach with \key{t} (session \texttt{main}). Closing Ghostty does not stop tmux; type
\key{t} again to get everything back.

\section{The status bar}

The bar sits at the top and is plain text:

\begin{plain}
 main  1 molly  2 humble      tunnel CPU 12% GPU 5% RAM 9.1/24G 48C D 1.2M U 120K docker 2 BAT 85%+  14:22
\end{plain}

\begin{keys}{3.4cm}
\head{Segment}{Meaning}
\texttt{main} & Session name. Blue normally, red while the prefix is pressed, grey when the outer tmux is turned off with \key{F12}\\
\texttt{1 molly} & Windows. \texttt{Z} zoomed pane, \texttt{*} activity in a background window, \texttt{S} panes typing in sync\\
\texttt{tunnel} & A \texttt{cloudflared} process is running (chapter~\ref{ch:web})\\
\texttt{CPU 12\%} & CPU load\\
\texttt{GPU 5\%} & GPU busy, from \texttt{ioreg} (no extra tool)\\
\texttt{RAM 9.1/24G} & Memory used / total\\
\texttt{48C} & CPU temperature, from \texttt{macmon} (Apple Silicon). Appears a few seconds after tmux starts\\
\texttt{D 1.2M U 120K} & Download / upload speed per second on the default network interface\\
\texttt{docker 2} & Running containers (hidden when Docker is not running)\\
\texttt{BAT 85\%+} & Battery; \texttt{+} means on power\\
\texttt{14:22} & Time\\
\end{keys}

The bar refreshes every 3 seconds. Every segment is a function in
\file{tmux/.config/tmux/scripts/status.sh}; run one by hand to debug it:
\cmd{~/.config/tmux/scripts/status.sh temp}.

\section{Keys}

\subsection{Panes and windows}

\begin{keys}{4.2cm}
\head{Keys}{Action}
\key{C-a c} & New window in the current folder\\
\key{C-a |} or \keyx{C-a \textbackslash} & Split right\\
\key{C-a -} & Split down\\
\key{C-h/j/k/l} & Move to the pane left/down/up/right; also moves between Neovim splits. In an SSH pane the key goes to the remote side\\
\key{C-a h/j/k/l} & Move pane, always the local tmux (also from SSH panes)\\
\key{C-a H/J/K/L} & Resize the pane (repeatable)\\
\key{C-a m} & Zoom the pane, again to restore\\
\key{C-a 1..9} & Go to window N\\
\key{C-a Tab} & Last window\\
\key{C-a n} / \key{C-a p} & Next / previous window\\
\key{C-a <} / \key{C-a >} & Move the window left / right\\
\key{C-a ,} & Rename the window\\
\key{C-a x} / \key{C-a X} / \key{C-a Q} & Close pane / window / session (asks first)\\
\key{C-a y} & Toggle typing into all panes of the window\\
\key{C-a v} & Copy mode (section~\ref{sec:clipboard})\\
\key{C-a C-l} & Clear the screen\\
\key{C-a r} & Reload the config\\
\end{keys}

\subsection{Popups and pickers}

\begin{keys}{2.4cm}
\head{Keys}{Action}
\key{C-a f} & Sessionizer: pick a project folder, open or switch to a session named after it\\
\key{C-a s} & Sessions and windows as a tree; \key{x} closes the selected one\\
\key{C-a S} & SSH hosts: \key{enter} new window, \key{ctrl-v} / \key{ctrl-s} split, \key{ctrl-o} here\\
\key{C-a d} & Running Docker containers: open a shell inside\\
\key{C-a a} & Web apps: start one in a new window (chapter~\ref{ch:web})\\
\key{C-a g} & lazygit in the current folder\\
\key{C-a D} & lazydocker\\
\key{C-a b} & btop\\
\key{C-a t} & Floating scratch terminal; its content stays when it is closed\\
\key{F12} & Turn the outer tmux off and on (for a tmux running on a robot)\\
\end{keys}

\screenshot[0.7]{ssh-picker}{\key{C-a S}: SSH host picker with the connection details of the
selected host.}

\section{Sessions per project}

The sessionizer (\key{C-a f}, or \cmd{ts} in the shell) lists folders under \file{~/apps},
\file{~/web}, \file{~/Projects}, \file{~/projects}, \file{~/dev}, \file{~/workspace},
\file{~/Developer}, \file{~/ros2_ws}, \file{~/dotfiles} and \file{~/Documents/GitHub}. Picking
one opens a session named after the folder, or switches to it when it exists.

\begin{keys}{3.6cm}
\head{Command}{Action}
\cmd{t} & Attach or create the session \texttt{main}\\
\cmd{ta}, \cmd{tl}, \cmd{tk <name>} & Attach, list, kill a session\\
\cmd{ts} & Sessionizer\\
\cmd{dev <host>} & Session named after an SSH host with windows \texttt{code}, \texttt{build}, \texttt{run}, each already connected (chapter~\ref{ch:containers})\\
\end{keys}

\section{Nested tmux}

When you SSH into a robot that runs its own tmux, the keys would go to the Mac's tmux. Press
\key{F12}: the session name turns grey, the bar shows \texttt{OFF(F12)} and every key goes to the
robot's tmux. Press \key{F12} again to come back. The robot's tmux keeps its default prefix
\key{C-b}, so the two never clash.
